% Independent four-interface cut proof candidate; not adopted.
\paragraph{Three safe-edge deletion cuts.}
Keep the opposite protected set $K$ fixed and put
$T=(K\cap E(B))\setminus P_B$. A Hamilton cycle in $B$ crosses
every nonempty proper vertex set an even positive number of times.
Thus if $B-T$ has just one edge across such a set $U$, every Hamilton
cycle in $B$ uses an edge of $T$, giving a protected common edge.
Only the Hamilton cycle on $B$ is quantified.

Use canonical target coordinates in direction F and the standard source
coordinates in direction R, always with the original interface first map.
The following $U$ are subsets of the full indicated component $B$.
\[
\begin{array}{c|c|c|c|c}
\text{source edit}&\text{direction}&B&U&\delta_{B-T}(U)\\\hline
00101:(3,3)&F&00121&x_0x_1=23,\ x_3\ne0&\{23203\!-23213\}\\
01014:(2,3)&F&01024&x_1x_2x_3=321&\{10212\!-13212\}\\
01024:(0,3),\ 01234:(0,2)&R,F&01024&x_0x_1=23&\{13130\!-23130\}
\end{array}
\]
The sets have six, five and ten vertices. Their full cuts have respectively
twelve, ten and fifteen other edges, all in $T$, by the legal one-coordinate
lists and the printed fixed protected sets and pools. The last two edits
give exactly the same $B-T$ because their original first colour maps are
inverse. No condition is imposed on the surviving edge's pool membership.
The global template construction still uses its specified pool and reserve
edges. This candidate replaces four individual proof premises only.
